\documentclass[10pt,a4paper]{amsart}
\usepackage{amsmath,amssymb,amsthm,graphicx}
\usepackage{verbatim}
\newenvironment{reference}{\comment}{\endcomment}
%\newenvironment{reference}{}{}
\newenvironment{slogan}{\comment}{\endcomment}
\newenvironment{history}{\comment}{\endcomment}
% For commutative diagrams we use Xy-pic
\usepackage[all]{xy}

% We use 2cell for 2-commutative diagrams.
\xyoption{2cell}
% Theorem environments.
%
\theoremstyle{plain}
\newtheorem{theorem}[subsection]{Theorem}
\newtheorem{proposition}[subsection]{Proposition}
\newtheorem{lemma}[subsection]{Lemma}

\theoremstyle{definition}
\newtheorem{definition}[subsection]{Definition}
\newtheorem{example}[subsection]{Example}
\newtheorem{exercise}[subsection]{Exercise}
\newtheorem{situation}[subsection]{Situation}

\theoremstyle{remark}
\newtheorem{remark}[subsection]{Remark}
\newtheorem{remarks}[subsection]{Remarks}

\numberwithin{equation}{subsection}

% Macros
%
\def\lim{\mathop{\mathrm{lim}}\nolimits}
\def\colim{\mathop{\mathrm{colim}}\nolimits}
\def\Spec{\mathop{\mathrm{Spec}}}
\def\Hom{\mathop{\mathrm{Hom}}\nolimits}
\def\Ext{\mathop{\mathrm{Ext}}\nolimits}
\def\SheafHom{\mathop{\mathcal{H}\!\mathit{om}}\nolimits}
\def\SheafExt{\mathop{\mathcal{E}\!\mathit{xt}}\nolimits}
\def\Sch{\mathit{Sch}}
\def\Mor{\mathop{\mathrm{Mor}}\nolimits}
\def\Ob{\mathop{\mathrm{Ob}}\nolimits}
\def\Sh{\mathop{\mathit{Sh}}\nolimits}
\def\NL{\mathop{N\!L}\nolimits}
\def\CH{\mathop{\mathrm{CH}}\nolimits}
\def\proetale{{pro\text{-}\acute{e}tale}}
\def\etale{{\acute{e}tale}}
\def\QCoh{\mathit{QCoh}}
\def\Ker{\mathop{\mathrm{Ker}}}
\def\Im{\mathop{\mathrm{Im}}}
\def\Coker{\mathop{\mathrm{Coker}}}
\def\Coim{\mathop{\mathrm{Coim}}}

% Boxtimes
%
\DeclareMathSymbol{\boxtimes}{\mathbin}{AMSa}{"02}

%
% Macros for moduli stacks/spaces
%
\def\QCohstack{\mathcal{QC}\!\mathit{oh}}
\def\Cohstack{\mathcal{C}\!\mathit{oh}}
\def\Spacesstack{\mathcal{S}\!\mathit{paces}}
\def\Quotfunctor{\mathrm{Quot}}
\def\Hilbfunctor{\mathrm{Hilb}}
\def\Curvesstack{\mathcal{C}\!\mathit{urves}}
\def\Polarizedstack{\mathcal{P}\!\mathit{olarized}}
\def\Complexesstack{\mathcal{C}\!\mathit{omplexes}}
% \Pic is the operator that assigns to X its picard group, usage \Pic(X)
% \Picardstack_{X/B} denotes the Picard stack of X over B
% \Picardfunctor_{X/B} denotes the Picard functor of X over B
\def\Pic{\mathop{\mathrm{Pic}}\nolimits}
\def\Picardstack{\mathcal{P}\!\mathit{ic}}
\def\Picardfunctor{\mathrm{Pic}}
\def\Deformationcategory{\mathcal{D}\!\mathit{ef}}
\usepackage{cleveref}
\crefname{originallemma}{original lemma}{original lemmas}
\Crefname{originallemma}{Original lemma}{Original lemmas}
\makeatletter
\newlabel{etale-lemma-characterize-closed-immersion}{{04XV}{source}{}{sourcepointer-1}{}}
\newlabel{etale-lemma-characterize-closed-immersion@cref}{{[originallemma][0][]04XV}{[1][0][]source}{}{}{}}
\makeatother
\crefname{originallemma}{original lemma}{original lemmas}
\Crefname{originallemma}{Original lemma}{Original lemmas}
\makeatletter
\newlabel{limits-lemma-Noetherian-dvr-valuative-proper}{{0208}{source}{}{sourcepointer-2}{}}
\newlabel{limits-lemma-Noetherian-dvr-valuative-proper@cref}{{[originallemma][0][]0208}{[1][0][]source}{}{}{}}
\makeatother
\crefname{originallemma}{original lemma}{original lemmas}
\Crefname{originallemma}{Original lemma}{Original lemmas}
\makeatletter
\newlabel{spaces-morphisms-lemma-base-change-finite-type}{{03XH}{source}{}{sourcepointer-3}{}}
\newlabel{spaces-morphisms-lemma-base-change-finite-type@cref}{{[originallemma][0][]03XH}{[1][0][]source}{}{}{}}
\makeatother
\crefname{originallemma}{original lemma}{original lemmas}
\Crefname{originallemma}{Original lemma}{Original lemmas}
\makeatletter
\newlabel{spaces-morphisms-lemma-closed-immersion-local}{{03M4}{source}{}{sourcepointer-4}{}}
\newlabel{spaces-morphisms-lemma-closed-immersion-local@cref}{{[originallemma][0][]03M4}{[1][0][]source}{}{}{}}
\makeatother
\crefname{originallemma}{original lemma}{original lemmas}
\Crefname{originallemma}{Original lemma}{Original lemmas}
\makeatletter
\newlabel{spaces-morphisms-lemma-composition-finite-type}{{03XG}{source}{}{sourcepointer-5}{}}
\newlabel{spaces-morphisms-lemma-composition-finite-type@cref}{{[originallemma][0][]03XG}{[1][0][]source}{}{}{}}
\makeatother
\crefname{originallemma}{original lemma}{original lemmas}
\Crefname{originallemma}{Original lemma}{Original lemmas}
\makeatletter
\newlabel{spaces-morphisms-lemma-locally-finite-type-locally-noetherian}{{04ZK}{source}{}{sourcepointer-6}{}}
\newlabel{spaces-morphisms-lemma-locally-finite-type-locally-noetherian@cref}{{[originallemma][0][]04ZK}{[1][0][]source}{}{}{}}
\makeatother
\crefname{originallemma}{original lemma}{original lemmas}
\Crefname{originallemma}{Original lemma}{Original lemmas}
\makeatletter
\newlabel{spaces-morphisms-lemma-properties-diagonal}{{03HK}{source}{}{sourcepointer-7}{}}
\newlabel{spaces-morphisms-lemma-properties-diagonal@cref}{{[originallemma][0][]03HK}{[1][0][]source}{}{}{}}
\makeatother
\title{Uniqueness of the lift along discretely valued diagrams implies separatedness of the morphism}
\author{The Stacks Project Authors}
\date{stacks-project:spaces-cohomology.tex:local-snapshot-20261012; GNU FDL 1.2 or later}
\begin{document}
\maketitle
\noindent\textit{Editorial note.} All delivered mathematical text is the author's own, copied from the pinned Stacks Project chapter file spaces-cohomology.tex; no statement or proof step was rewritten, reordered or removed.\par\medskip
\noindent\textit{Editorial note.} Preamble provenance: the pinned chapter file opens with \textbackslash{}input\{preamble\} and carries no preamble text of its own, so the project's single Stacks-wide preamble staged at sources/stacks-spaces-cohomology/preamble.tex is excerpted verbatim as three spans, lines 8-12 (verbatim package and the project's own reference/slogan/history display environments), lines 14-18 (Xy-pic setup) and lines 154-217 (theorem declarations and mathematical macros). The wrapper class and the amsmath/amssymb/amsthm/graphicx base packages are supplied by the assembler. No macro was invented or redefined.\par\medskip
\noindent\textit{Editorial note.} Same-chapter pointers are located at their \textbackslash{}label lines inside the pinned chapter file; cross-chapter pointers are located in the project's own tag list staged in the snapshot (sources/stacks-chow/tags.txt), at the verbatim lines carrying their canonical tags.\par\medskip
\noindent\textit{Allowed standard prerequisites (editorial).}\begin{itemize}
\item The properties of the diagonal of a morphism of algebraic spaces. Morphisms of Spaces, label spaces-morphisms-lemma-properties-diagonal (invoked at proof line 3655). It gives that \$\textbackslash{}Delta\$ is representable, separated, a monomorphism and locally of finite type, which is the starting point; representability makes the fibre product \$V\$ below a scheme.
\item Closed immersions are detected etale locally on the target. Morphisms of Spaces, label spaces-morphisms-lemma-closed-immersion-local (invoked at proof line 3660). It reduces the closedness of \$\textbackslash{}Delta\$ to the closedness of the map \$V \textbackslash{}to U\$ obtained from an affine etale \$U \textbackslash{}to X \textbackslash{}times\_Y X\$.
\item The resulting auxiliary scheme is Noetherian, of finite type and quasi-compact. Morphisms of Spaces, labels spaces-morphisms-lemma-composition-finite-type, spaces-morphisms-lemma-base-change-finite-type and spaces-morphisms-lemma-locally-finite-type-locally-noetherian (invoked at proof lines 3662-3665). They make \$U\$ Noetherian and, together with quasi-separatedness of \$f\$, show that \$V \textbackslash{}to U\$ is a quasi-compact morphism of finite type, so that the valuative criterion applies.
\item The valuative criterion of properness for a morphism of Noetherian schemes. Limits, label limits-lemma-Noetherian-dvr-valuative-proper (invoked at proof line 3681). Once the uniqueness of the dotted arrow is verified for all discrete valuation ring diagrams over \$U\$, it gives that \$V \textbackslash{}to U\$ is proper.
\item A proper monomorphism is a closed immersion. Etale Morphisms, label etale-lemma-characterize-closed-immersion (invoked at proof line 3684). Since \$\textbackslash{}Delta\$ is a monomorphism and has been shown proper, it is a closed immersion, i.e. \$f\$ is separated.
\end{itemize}\medskip
\ifdefined\hypertarget\hypertarget{sourcepointer-1}{}\fi
\noindent\textit{Source reference.} Original lemma 04XV: sources/stacks-chow/tags.txt, source lines 6408--6408. Navigation pointer for \textbackslash{}ref\{etale-lemma-characterize-closed-immersion\} used inside the counted span (line 3685); the argument that consumes it is retained in the span.\par\medskip
\ifdefined\hypertarget\hypertarget{sourcepointer-2}{}\fi
\noindent\textit{Source reference.} Original lemma 0208: sources/stacks-chow/tags.txt, source lines 2625--2625. Navigation pointer for \textbackslash{}ref\{limits-lemma-Noetherian-dvr-valuative-proper\} used inside the counted span (line 3681); the argument that consumes it is retained in the span.\par\medskip
\ifdefined\hypertarget\hypertarget{sourcepointer-3}{}\fi
\noindent\textit{Source reference.} Original lemma 03XH: sources/stacks-chow/tags.txt, source lines 5148--5148. Navigation pointer for \textbackslash{}ref\{spaces-morphisms-lemma-base-change-finite-type\} used inside the counted span (line 3664); the argument that consumes it is retained in the span.\par\medskip
\ifdefined\hypertarget\hypertarget{sourcepointer-4}{}\fi
\noindent\textit{Source reference.} Original lemma 03M4: sources/stacks-chow/tags.txt, source lines 4736--4736. Navigation pointer for \textbackslash{}ref\{spaces-morphisms-lemma-closed-immersion-local\} used inside the counted span (line 3660); the argument that consumes it is retained in the span.\par\medskip
\ifdefined\hypertarget\hypertarget{sourcepointer-5}{}\fi
\noindent\textit{Source reference.} Original lemma 03XG: sources/stacks-chow/tags.txt, source lines 5147--5147. Navigation pointer for \textbackslash{}ref\{spaces-morphisms-lemma-composition-finite-type\} used inside the counted span (line 3663); the argument that consumes it is retained in the span.\par\medskip
\ifdefined\hypertarget\hypertarget{sourcepointer-6}{}\fi
\noindent\textit{Source reference.} Original lemma 04ZK: sources/stacks-chow/tags.txt, source lines 6468--6468. Navigation pointer for \textbackslash{}ref\{spaces-morphisms-lemma-locally-finite-type-locally-noetherian\} used inside the counted span (line 3665); the argument that consumes it is retained in the span.\par\medskip
\ifdefined\hypertarget\hypertarget{sourcepointer-7}{}\fi
\noindent\textit{Source reference.} Original lemma 03HK: sources/stacks-chow/tags.txt, source lines 4570--4570. Navigation pointer for \textbackslash{}ref\{spaces-morphisms-lemma-properties-diagonal\} used inside the counted span (line 3655); the argument that consumes it is retained in the span.\par\medskip
\begin{lemma}
\label{lemma-check-separated-dvr}
Let $S$ be a scheme. Let $f : X \to Y$ be a morphism of algebraic
spaces over $S$. Assume
\begin{enumerate}
\item $Y$ is locally Noetherian,
\item $f$ is locally of finite type and quasi-separated,
\item for every commutative diagram
$$
\xymatrix{
\Spec(K) \ar[r] \ar[d] & X \ar[d] \\
\Spec(A) \ar[r] \ar@{-->}[ru] & Y
}
$$
where $A$ is a discrete valuation ring and $K$ its fraction field,
there is at most one dotted arrow making the diagram commute.
\end{enumerate}
Then $f$ is separated.
\end{lemma}

\begin{proof}
We have to show that the diagonal $\Delta : X \to X \times_Y X$ is a closed
immersion. We already know $\Delta$ is representable, separated,
a monomorphism, and locally of finite type, see Morphisms of Spaces, Lemma
\ref{spaces-morphisms-lemma-properties-diagonal}.
Choose an affine scheme $U$ and an \'etale morphism
$U \to X \times_Y X$. Set $V = X \times_{\Delta, X \times_Y X} U$.
It suffices to show that $V \to U$ is a closed immersion
(Morphisms of Spaces, Lemma
\ref{spaces-morphisms-lemma-closed-immersion-local}).
Since $X \times_Y X$ is locally of finite type over $Y$ we see that
$U$ is Noetherian (use Morphisms of Spaces, Lemmas
\ref{spaces-morphisms-lemma-composition-finite-type},
\ref{spaces-morphisms-lemma-base-change-finite-type}, and
\ref{spaces-morphisms-lemma-locally-finite-type-locally-noetherian}).
Note that $V$ is a scheme as $\Delta$ is representable.
Also, $V$ is quasi-compact because $f$ is quasi-separated.
Hence $V \to U$ is of finite type. Consider a commutative diagram
$$
\xymatrix{
\Spec(K) \ar[r] \ar[d] & V \ar[d] \\
\Spec(A) \ar[r] \ar@{-->}[ru] & U
}
$$
of morphisms of schemes
where $A$ is a discrete valuation ring with fraction field $K$.
We can interpret the composition $\Spec(A) \to U \to X \times_Y X$
as a pair of morphisms $a, b : \Spec(A) \to X$ agreeing as morphisms
into $Y$ and equal when restricted to $\Spec(K)$. Hence our assumption
(3) guarantees $a = b$ and we find the dotted arrow in the diagram.
By Limits, Lemma \ref{limits-lemma-Noetherian-dvr-valuative-proper}
we conclude that $V \to U$ is proper. In other words, $\Delta$ is proper.
Since $\Delta$ is a monomorphism, we find that $\Delta$ is a
closed immersion (\'Etale Morphisms, Lemma
\ref{etale-lemma-characterize-closed-immersion}) as desired.
\end{proof}

\end{document}
